\documentclass[c,18pt]{beamer}
\listfiles

\mode<presentation>
{
  \usetheme[deutsch,titlepage0]{KIT}
  \setbeamercovered{transparent}
  \setbeamertemplate{enumerate items}[ball]
}
\usepackage[utf8]{inputenc}
\date{20.11.2018}

\newlength{\Ku}
\setlength{\Ku}{1.43375pt}

\usepackage{templateSlide/exercises}
\usepackage[java]{code}

\makeatletter
\def\input@path{{uebungsfolien/}}
\graphicspath{{uebungsfolien/}}
\makeatother

%% title slide
\title[Übung 05: Einstieg in die Informatik und algorithmische Mathematik]
  {Übung 05: Einstieg in die Informatik und algorithmische Mathematik}
\subtitle{Albert Mink | Wintersemester 2018/2019}

\author[Albert Mink, ]{KIT}

\AuthorTitleSep{\relax}

\institute[Institut für Angewandte und Numerische Mathematik (IANM)]{Institut für Angewandte und Numerische Mathematik}

\TitleImage[width=\titleimagewd]{logos/KIT-Titel}
\logo{\includegraphics{logos/lbrg_logo}}

%%%%%%%%%%%%%%%%%%%%%%%%%%%%%%%%%%%%%%%%%%%%%%%%%%%%%%%%%%%%%%%%%%%%%%%%
% Document
%%%%%%%%%%%%%%%%%%%%%%%%%%%%%%%%%%%%%%%%%%%%%%%%%%%%%%%%%%%%%%%%%%%%%%%%
\begin{document}
\begin{frame}
  \maketitle
\end{frame}

%%%%%%%%%%%%%%%%%%%%%%%%%%%%%%%%%%%%%%%%%%%%%%%%%%%%%%%%%%%%%%%%%%%%%%%%

%%%%%%%%%%%%%%%%%%%%%%%%%%%%%%%%%%%%%%%%%%%%%%%%%%%%%%%%%%%%%%%%%%%%%%%%
%%%%%%%%%%%%%%%%%%%%%%%%%%%%%%%%%%%%%%%%%%%%%%%%%%%%%%%%%%%%%%%%%%%%%%%%
\begin{frame}
  \frametitle{Arbeitsblatt 5}%
\tableofcontents
\end{frame}
\setcounter{exercise}{20}
%%%%%%%%%%%%%%%%%%%%%%%%%%%%%%%%%%%%%%%%%%%%%%%%%%%%%%%%%%%%%%%%%%%%%%%%


%%%%%%%%%%%%%%%%%%%%%%%%%%%%%%%%%%%%%%%%%%%%%%%%%%%%%%%%%%%%%%%%%%%%%%%%
%%%%%%%%%%%%%%%%%%%%%%%%%%%%%%%%%%%%%%%%%%%%%%%%%%%%%%%%%%%%%%%%%%%%%%%%
\section{Wiederholung}\label{K:wdh}
\begin{frame}
  \frametitle{\ref{K:wdh} Wiederholung}%
\tableofcontents[current]
\end{frame}
%%%%%%%%%%%%%%%%%%%%%%%%%%%%%%%%%%%%%%%%%%%%%%%%%%%%%%%%%%%%%%%%%%%%%%%%

%%%%%%%%%%%%%%%%%%%%%%%%%%%%%%%%%%%%%%%%%%%%%%%%%%%%%%%%%%%%%%%%%%%%%%%%
%%%%%%%%%%%%%%%%%%%%%%%%%%%%%%%%%%%%%%%%%%%%%%%%%%%%%%%%%%%%%%%%%%%%%%%%
\def\stitle{Erzeugen von Feldern}
\subsection{\stitle}\label{S:Erzeugen}
\begin{frame}[t]%
  \frametitle{\ref{K:wdh}.\ref{S:Erzeugen} \stitle}
\medskip

Ein Feld kann mehrere Variablen vom selben Datentyp enthalten, hier \code{int}.
\lstinputlisting[style=JAVAlines]{wdh-felder/FelderErzeugen.java}

\end{frame}
%%%%%%%%%%%%%%%%%%%%%%%%%%%%%%%%%%%%%%%%%%%%%%%%%%%%%%%%%%%%%%%%%%%%%%%%


%%%%%%%%%%%%%%%%%%%%%%%%%%%%%%%%%%%%%%%%%%%%%%%%%%%%%%%%%%%%%%%%%%%%%%%%
%%%%%%%%%%%%%%%%%%%%%%%%%%%%%%%%%%%%%%%%%%%%%%%%%%%%%%%%%%%%%%%%%%%%%%%%
\def\stitle{Zugreifen auf Felder}
\subsection{\stitle}\label{S:Zugreifen}
\begin{frame}[t]%
  \frametitle{\ref{K:wdh}.\ref{S:Zugreifen} \stitle}
\medskip

Mit dem Operator \code{[]} wird auf bestimmte Feldkomponenten zugegriffen.
\lstinputlisting[style=JAVAlines]{wdh-felder/FelderZugreifen.java}

\end{frame}
%%%%%%%%%%%%%%%%%%%%%%%%%%%%%%%%%%%%%%%%%%%%%%%%%%%%%%%%%%%%%%%%%%%%%%%%

%%%%%%%%%%%%%%%%%%%%%%%%%%%%%%%%%%%%%%%%%%%%%%%%%%%%%%%%%%%%%%%%%%%%%%%%
%%%%%%%%%%%%%%%%%%%%%%%%%%%%%%%%%%%%%%%%%%%%%%%%%%%%%%%%%%%%%%%%%%%%%%%%
\def\stitle{Felder Bibliothek in Java}
\subsection{\stitle}\label{S:bequem}
\begin{frame}[t]%
  \frametitle{\ref{K:wdh}.\ref{S:bequem} \stitle}
\medskip

In \code{java.util.Arrays} werden n\"utzliche Funktion der Java Bibliothek bereit gestellt.
\lstinputlisting[style=JAVAlines]{wdh-felder/FelderAdv.java}

\end{frame}
%%%%%%%%%%%%%%%%%%%%%%%%%%%%%%%%%%%%%%%%%%%%%%%%%%%%%%%%%%%%%%%%%%%%%%%%


%%%%%%%%%%%%%%%%%%%%%%%%%%%%%%%%%%%%%%%%%%%%%%%%%%%%%%%%%%%%%%%%%%%%%%%%
%%%%%%%%%%%%%%%%%%%%%%%%%%%%%%%%%%%%%%%%%%%%%%%%%%%%%%%%%%%%%%%%%%%%%%%%
\def\stitle{Felder Beispiel Anwendung, Ger\"ust}
\subsection{\stitle}\label{S:BeispielG}
\begin{frame}[t]%
  \frametitle{\ref{K:wdh}.\ref{S:BeispielG} \stitle}
\medskip

Lese Vektor ein und berechne die Norm.
\lstinputlisting[style=JAVAsmalllines]{wdh-felder/FelderAnwBare.java}

\end{frame}
%%%%%%%%%%%%%%%%%%%%%%%%%%%%%%%%%%%%%%%%%%%%%%%%%%%%%%%%%%%%%%%%%%%%%%%%


%%%%%%%%%%%%%%%%%%%%%%%%%%%%%%%%%%%%%%%%%%%%%%%%%%%%%%%%%%%%%%%%%%%%%%%%
%%%%%%%%%%%%%%%%%%%%%%%%%%%%%%%%%%%%%%%%%%%%%%%%%%%%%%%%%%%%%%%%%%%%%%%%
\def\stitle{Felder Beispiel Anwendung}
\subsection{\stitle}\label{S:Beispiel}
\begin{frame}[t]%
  \frametitle{\ref{K:wdh}.\ref{S:Beispiel} \stitle}
\medskip

Lese Vektor ein und berechne die Norm.
\lstinputlisting[style=JAVAsmalllines]{wdh-felder/FelderAnw.java}

\end{frame}
%%%%%%%%%%%%%%%%%%%%%%%%%%%%%%%%%%%%%%%%%%%%%%%%%%%%%%%%%%%%%%%%%%%%%%%%


\setcounter{exercise}{17}
\begin{frame}[t]%
\medskip

\begin{exercise}{Felder: Grundlagen}
\begin{body}
Felder k"onnen aus den Datentypen \code{double}, \code{int}, \code{char}, \code{String}, $\ldots$ bestehen.
Der Index aller Felder beginnt mit '0', d.h. ein Feld mit 10 Elementen hat die Indizes von 0 bis 9.
\medskip

\lstinputlisting[style=JAVAlines]{felder-grundl/FelderGrundlagen.java}

Was wird auf dem Bildschirm ausgegeben?
\end{body}
\end{exercise}
\end{frame}

\setcounter{exercise}{18}
\def\stitle{\theexercise\ - Felder: Binäre Suche}
\section{\stitle}

\begin{frame}[t]%
    \frametitle{\stitle}
\medskip
Gegeben sei untenstehende Implementierung der Binären Suche.
\lstinputlisting[style=JAVAsmalllines]{\getexercisefolder/Binsearch.java}
\end{frame}


\begin{frame}
 \frametitle{\stitle \, - Aufgabenstellung}%
\medskip

Hier ein Programmausschnitt.
Welche Ausgabe liefert das Programm f"ur die Eingaben 9 bzw. 10 zur"uck?
\lstinputlisting[style=JAVAlines]{\getexercisefolder/Binsearch_snippet.java}
\end{frame}

\setcounter{exercise}{19}
\begin{frame}[t]%
  \setcounter{exercise}{22} % TODO make exercise number dynamic
  \medskip
  \begin{exercise}{Felder (a)}
  \begin{body}
  Betrachten Sie den folgenden Ausschnitt eines Java-Programms.
  \lstinputlisting[style=JAVAlines]{felder-felder/Feldera.java}
  \begin{parts}
  \item Was wird auf der Konsole ausgegeben?
  \pause
  \item \code{[1, 3, 6, 10, 15]}
  \end{parts}
  \end{body}
  \end{exercise}
\end{frame}


\begin{frame}[t]%
  \setcounter{exercise}{22} % TODO make exercise number dynamic
  \medskip
  \begin{exercise}{Felder (b)}
  \begin{body}
  Betrachten Sie den folgenden Ausschnitt eines Java-Programms.
  \lstinputlisting[style=JAVAlines]{felder-felder/Felderb.java}
  \begin{parts}
  \item Was wird auf der Konsole ausgegeben?
  \pause
  \item \code{[1, 3, 5, 7, 9]}
  \end{parts}
  \end{body}
  \end{exercise}
\end{frame}

\begin{frame}[t]%
  \setcounter{exercise}{22} % TODO make exercise number dynamic
  \medskip
  \begin{exercise}{Felder (c)}
  \begin{body}
  Betrachten Sie den folgenden Ausschnitt eines Java-Programms.
  \lstinputlisting[style=JAVAlines]{felder-felder/Felderc.java}
  \begin{parts}
  \item Was wird auf der Konsole ausgegeben?
  \pause
  \item \code{[1, 2, 3, 2, 1]}
  \end{parts}
  \end{body}
  \end{exercise}
\end{frame}


\begin{frame}[t]%
  \setcounter{exercise}{22} % TODO make exercise number dynamic
  \medskip
  \begin{exercise}{Felder (d)}
  \begin{body}
  Betrachten Sie den folgenden Ausschnitt eines Java-Programms.
  \lstinputlisting[style=JAVAlines]{felder-felder/Felderd.java}
  \begin{parts}
  \item Was wird auf der Konsole ausgegeben?
  \pause
  \item \code{[0, 0, 3, 3, 3]}
  \end{parts}
  \end{body}
  \end{exercise}
\end{frame}

\begin{frame}[t]%
  \setcounter{exercise}{22} % TODO make exercise number dynamic
  \medskip
  \begin{exercise}{Felder (e)}
  \begin{body}
  Betrachten Sie den folgenden Ausschnitt eines Java-Programms.
  \lstinputlisting[style=JAVAlines]{felder-felder/Feldere.java}
  \begin{parts}
  \item Was wird auf der Konsole ausgegeben?
  \pause
  \item \code{[1, 4, 9, 16, 25]}
  \end{parts}
  \end{body}
  \end{exercise}
\end{frame}

\begin{frame}[t]%
\setcounter{exercise}{22} % TODO make exercise number dynamic
  \medskip
  \begin{exercise}{Felder (f)}
  \begin{body}
  Betrachten Sie den folgenden Ausschnitt eines Java-Programms.
  \lstinputlisting[style=JAVAlines]{felder-felder/Felderf.java}
  \begin{parts}
  \item Was wird auf der Konsole ausgegeben?
  \pause
  \item \code{[5, 4, 3, 2, 1]}
  \end{parts}
  \end{body}
  \end{exercise}
\end{frame}


\begin{frame}
\centering
\Huge\textcolor{KITgreen}{Fragen?}
\vspace{2cm}

{\LARGE
N\"achste \"Ubung: 27. November\\
Besprechung Arbeitsblatt 6
}
\end{frame}


%%%%%%%%%%%%%%%%%%%%%%%%%%%%%%%%%%%%%%%%%%%%%%%%%%%%%%%%%%%%%%%%%%%%%%%%
\end{document}
